\documentclass{aci}

%%%%%%%%%%%%%%%%%%%%%%%%%%%%%%%%%%%%%%%%%%
\usepackage{txfonts}
\usepackage{booktabs}
\usepackage{longtable}
\usepackage{lipsum}
\usepackage{hyperref}
\hypersetup{colorlinks=true}
\usepackage{subcaption}
\usepackage{placeins}
\usepackage{tikz}
\usetikzlibrary{automata,positioning,arrows}

%%%%%%%%%%%%%%%%%%%%%%%%%%%%%%%%%%%%%%%%%%
\usepackage[backend=bibtex,sorting=none]{biblatex}
\addbibresource{./references/paper.bib}
\addbibresource{./references/unsorted.bib}
\addbibresource{./references/def_model.bib}

%%%%%%%%%%%%%%%%%%%%%%%%%%%%%%%%%%%%%%%%%%
\newcommand{\ep}{\varepsilon}
\newcommand{\eps}[1]{{#1}_{\varepsilon}}

\def\typeofarticle{Research Article} \def\currentvolume{1} \def\currentissue{1}
\def\currentyear{2021} \def\currentmonth{} \def\ppages{1--30} \def\DOI{to
appear} \def\Received{June 2022} \def\Accepted{July 2022} \def\Published{October
2022 }
\numberwithin{equation}{section}
\DeclareMathOperator*{\essinf}{ess\,inf}

\usepackage{array}
\newcolumntype{L}{>{\arraybackslash}m{8cm}}

%%%%%%%%%%%%%%%%%%%%%%%%%%%%%%%%%%%%%%%%%%
\begin{document}

\title{Survey on Definition Modeling}

\author{%
    Noah Gardner\affil{1},
    Hafiz Khan\affil{2},
    Chih-Cheng Hung\affil{1}\corrauth
}%

% \shortauthors is used in copyright information in the end of the paper
\shortauthors{the Author(s)}

\address{%
    \addr{\affilnum{1}}{Laboratory of Machine Vision and Security Research, %
        College of Computing and Software Engineering, Kennesaw State %
        University, Marietta GA, USA}%
    \addr{\affilnum{2}}{Laboratory of Ubiquitous Data Mining, %
        College of Computing and Software Engineering, Kennesaw State %
        University, Marietta GA, USA}}
\corraddr{chung1@kennesaw.edu}

\begin{abstract}
    Definition modeling, the task of generating a definition for a given term,
    is a relatively new area of research applied in the evaluation of word
    embeddings that is gathering attention. Automatic generation of dictionary
    quality definitions has many applications in natural language processing,
    such as sentiment analysis, machine translation, and word sense
    disambiguation. Additionally, definition modeling is also useful for
    evaluating the quality of word embeddings. As more research is done in this
    field, the need for a summary of different applications, approaches, and
    obstacles grows apparent. This survey aims to provide an overview of the
    current research in definition modeling, and to provide a list of future
    directions and trends.
\end{abstract}
\keywords{Definition modeling, definition generation}
\maketitle

\section{Introduction}
Definitions are explicit representations of words or phrases that are valuable
for exposing the aspects of a given term. In general, definitions are
unambiguous and succinct: they should be easy to read and understand. Recent research
has allowed the creation of neural language models that can generate useful
definitions from embeddings \cite{bosc_auto_2018, hill_learning_2016, washio_bridging_2019}. Word embeddings are vector representations of words that have been employed in a variety of \textit{natural language
    processing} (NLP) tasks. They are useful for capturing lexical syntax and
semantic similarity. Mikolov et al. \cite{mikolov_distributed_2013} have shown that
basic mathematical operations applied to word embeddings can have meaningful
language understanding. However, as continuous representations, the
interpretability of word embeddings is limited.

\begin{figure}[h]
    \centering
    \includegraphics[width=.9\textwidth]{assets/figures/monoseme.png}
    \caption{Monoseme example word and definition. A definition model could
        generate the definition \textit{producing a great deal of profit} for the
        input word \textit{lucrative}.}
    \label{fig:monoseme}
\end{figure}

The problem of \textit{definition modeling} was proposed by Noraset et al.~\cite{noraset_definition_2016} to evaluate word embeddings. The task of definition modeling is to generate a definition for a given term. The goal of a model trained on this task is to train on word embedding and definition pairs to learn to generate a definition for a given word or phrase. An example of a \textit{monoseme} (word with a single definition) is given in Figure~\ref{fig:monoseme}. Given the input word \textit{lucrative}, a model trained on
the task of definition modeling would produce the output definition
\textit{producing a great deal of profit}.

\begin{figure}[h]
    \centering
    \includegraphics[width=.9\textwidth]{assets/figures/polyseme.png}
    \caption{Polyseme example word and definition. A definition model could generate any of the target definitions.}
    \label{fig:polyseme}
\end{figure}

In addition to being a relatively new language modeling task, definition
modeling has attracted attention from the literature in a number of areas.
First, it was shown that the definition model has poor performance when
generating definitions for polysemes: words with multiple definitions
\cite{gadetsky_conditional_2018}. An example polysemous word is shown in Figure
\ref{fig:polyseme}. Given the input \textit{word}, the goal of a definition
model would be to generate one of the target definitions, most ideally the
closest definition to the word sense of the input. However, it is challenging to
know the word sense given only the input word.

The problem of polysemous words was not addressed in the original work, as only one definition mapped to each word. Once researchers attempted to address this problem, they found that the definition model could not learn the semantics of the polyseme with only the word as an input. Therefore, it was necessary to augment the definition model with additional information, namely, an example sentence that sets the word to be defined (\textit{definiendum}) inside to
provide context. This method has been shown to alleviate the problem of generating definitions for polysemes and also improve the performance of the definition model on several measures~\cite{bevilacqua_generationary_2020,
    gadetsky_conditional_2018, mickus_mark_2019}.

\begin{figure}[h]
    \centering
    \includegraphics[width=.9\textwidth]{assets/figures/context.png}
    \caption{Context-dependent definition task. The word to be defined is marked
        in bold. In this case, although \textbf{word} is a polyseme, there is only
        one correct target definition due to the contextual information provided.}
    \label{fig:context_poly}
\end{figure}

Definition modeling, especially as a sequence-to-sequence task, is similar to other NLP tasks, such as \textit{word-sense disambiguation}, \textit{word-in-context}, and \textit{definition extraction}~\cite{bevilacqua_generationary_2020, huang_cdm_2021}. When using context to generate a definition from an input word, the input's word sense must be extracted to select the correct definition. The goal of word-sense disambiguation is similar in that the goal of word-sense disambiguation is to identify the sense of a word used in a sentence. Similarly, definition extraction seeks to extract definitions of terms from an existing corpus~\cite{huang_cdm_2021}. Figure~\ref{fig:context_poly} shows an example of definition modeling in a context-dependent situation. In the example, a reference context is given. Inside the reference context, a target word
\textit{word} is marked as the word to be defined. The goal of a definition model given this context and marked word would be to generate the target definition \textit{a command, password, or signal}.

Our paper is organized into three sections. Section 2 reviews definition modeling methods as well as word embeddings. Section 3 shares benchmark datasets
and statistics that can be used when formulating and evaluating a definition modeling method. Section 4 explores challenges encountered in this research field and gives suggestions for future work.

\section{Related Work}

\section{Scope of this Review}
Scope of this Review: breif overview of the review categories: examples - tasks,
data types, methods, future directions and challenges.

\section{Methodologies Explored}
main key points -- definition generation and word embedding (similarity disabmiguity etc. polysemy). Few line of text should be here.  

\begin{longtable}{|p{3.5cm}|p{3.5cm}|p{3.5cm}|p{3.5cm}|}
    \caption{Definition Modeling Methods}                                                                                                                                                                                                                                                                                                                                                   \\
    \hline
    Article                                                                                                                   & Evaluation                                                       & Dataset                                                                                                                                       & Models                                   \\
    \hline
    \citeauthor*{noraset_definition_2016} \citeyear{noraset_definition_2016} \cite{noraset_definition_2016}                   & Perplexity, BLEU                                                 & GCIDE, Custom                                                                                                                                 & RNN, Word2Vec, LSTM                      \\
    \hline
    \citeauthor*{gadetsky_conditional_2018} \citeyear{gadetsky_conditional_2018} \cite{gadetsky_conditional_2018}             & Perplexity, BLEU                                                 & Oxford                                                                                                                                        & LSTM, Word2Vec, SkipGram                 \\
    \hline
    \citeauthor*{chang_what_2019} \citeyear{chang_what_2019} \cite{chang_what_2019}                                           & Precision, ROUGE-L, Cosine similarity                            & Oxford                                                                                                                                        & ELMo, BERT, FastText                     \\
    \hline
    \citeauthor*{washio_bridging_2019} \citeyear{washio_bridging_2019} \cite{washio_bridging_2019}                            & Perplexity, BLEU                                                 & \cite{noraset_definition_2016}, \cite{gadetsky_conditional_2018}                                                                              & Encoder/decoder                          \\
    \hline
    \citeauthor*{mickus_mark_2019} \citeyear{mickus_mark_2019} \cite{mickus_mark_2019}                                        & Perplexity                                                       & \cite{noraset_definition_2016}, \cite{gadetsky_conditional_2018}, Custom                                                                      & GloVe, Transformer, sequence-to-sequence \\
    \hline
    \citeauthor*{li_explicit_2020} \citeyear{li_explicit_2020} \cite{li_explicit_2020}                                        & BLEU, METEOR, Human                                              & WordNet, Oxford                                                                                                                               & CNN, LSTM                                \\
    \hline
    \citeauthor*{bevilacqua_generationary_2020} \citeyear{bevilacqua_generationary_2020} \cite{bevilacqua_generationary_2020} & Perplexity, BLEU, ROUGE-L, METEOR, BERTScore                     & Oxford \cite{chang_what_2019}, Sem-Cor \cite{miller_semantic_1993}, Wiktionary, GCIDE \cite{noraset_definition_2016}, Hei++ (Custom), WordNet & BART                                     \\
    \hline
    \citeauthor*{sojka_evaluating_2020} \citeyear{sojka_evaluating_2020} \cite{sojka_evaluating_2020}                         & BLEU, rBLEU (recall-based), fBLEU (harmonic mean of BLEU, rBLEU) & Wiktionary, OmegaWiki, WordNet                                                                                                                & AdaGram, Word2Vec, CNN, RNN              \\
    \hline
    \citeauthor*{huang_cdm_2021} \citeyear{huang_cdm_2021} \cite{huang_cdm_2021}                                              & BLEU, ROUGE-L, METEOR, BERTScore                                 & Wikipedia                                                                                                                                     & BERT                                     \\
    \hline
    \citeauthor*{reid_vcdm_2020} \citeyear{reid_vcdm_2020} \cite{reid_vcdm_2020}                                              & BLEU, BERTScore (Precision, Recall, F1), Human                   & Oxford, Urban, Wikipedia, Cambridge, Robert (French)                                                                                          & BERT, LSTM                               \\
    \hline
    % \citeauthor*{noraset_definition_2016} \citeyear{noraset_definition_2016} \cite{noraset_definition_2016}                   & Perplexity, BLEU                             & GCIDE, Custom                                                                                                                                 & RNN, Word2Vec, LSTM                      \\
\end{longtable}

\subsection{Language Models}
A definition model is a language model that is trained on a set of definitions
\cite{noraset_definition_2016}. The goal of a definition model is to learn to
generate a definition ($\textbf{d} = [d_1, ..., d_T]$) for a given term $w$. The
probability of generating the $t$-th word in a definition depends on both the
previous words in the definition and the word to be defined (Eq.
\ref{eq:definition_model}).

\begin{equation}
    \label{eq:definition_model}
    p(\textbf{d} | w) = \prod_{t=1}^{T} p(d_t | d_1,...,d_{t-1}, w)
\end{equation}

Noraset et al. condition an RNN to generate a defintion from an input seed word.
They modify the model by updating the output of the recurrent unit with an
update function inspired by GRU update gate \cite{noraset_definition_2016}. They
apply pretrained word embeddings generated from Word2Vec.

In later work, it was shown that the definition model does not generate
definitions for words with ambigious word sense, especially polysemantic words.
The following context-aware definition model was proposed by
\citeauthor*{gadetsky_conditional_2018} in order to tackle this challenge
\cite{gadetsky_conditional_2018}. They extend Equation \ref{eq:definition_model}
by adding a context term ($\textbf{c} = [c_1, ..., c_T]$) which is a context or
example sentence to be used in the generation of the definition.

\begin{equation}
    \label{eq:context_aware_definition_model}
    p(\textbf{d} | w, \textbf{c}) = \prod_{t=1}^{T} p(d_t | d_1,...,d_{t-1}, w, \textbf{c})
\end{equation}

In order to generate a definition, they use an attention-based SkipGram model to
extract dimensions from the embedding which contain the most relevant
information.

\cite{huang_cdm_2021} focused on language generation, etc.\cite{gadetsky_conditional_2018} conditional generation of definition.\cite{zheng_decompose_2021}, chienese defition generation.
\cite{sojka_evaluating_2020}, generation
\cite{li_explicit_2020}
\cite{bevilacqua_generationary_2020}
\cite{zhang_improving_2020}
\cite{yang_incorporating_2020}
\cite{ishiwatari_learning_2019}a general task of defining un-
known phrases given their contexts.

\subsection{Word Embedding}
In recent years, definition modeling has gained popularity, and researchers proposed various methods to map between definition and the target word. One of the significant issues in mapping words with the dictionary is contextual ambiguity and embedding-based word representation. In NLP, word embedding represents the vector representation of words that encodes the word's meaning such that similar words have similar vector space. There are various word embedding techniques used to resolve ambiguity between words. In the definition modeling problem, word to vector representation is a key factor in modeling definition words/terms for a given term. 

\cite{bosc_auto_2018} exploited dictionary recursivity into consideration and proposed an autoencoder-based word embedding algorithm, and generated a single embedding per word—the proposed auto-encoder model comprises of LSTM encoder and decoder. The authors introduced three embeddings - i) definition embeddings produced by the proposed definition encoder, ii) input embeddings for the encoder, and iii) output embeddings. While modeling these embedding models, the author incorporates consistency penalty as soft weight in their cost function to enforce input embedding and definition embedding closer. 

In \cite{washio_bridging_2019}, the authors consider lexical-semantic relations between the defined word and defining words using unsupervised methods to propose definition modeling. To learn embedding author proposed LSTM based encoder and decoder with additional cost function to learn word-pair embeddings in the decoder and capture lexical-semantic relations.

Dictionary embeddings often follow a genus + differentia structure for a dictionary definition~\cite{noraset_definition_2016}—author capture hypernyms embedding following proper genus database WebIsA containing hypernym relations. In addition, the author incorporates char-CNN to capture affixes to model gated-RNN based definition modeling.


Word-embeddings are learned from large corpora. Therefore, it may consists of biases such as gender, race, religion, etc. On the otherhand, word dictionaries contain unbiased, concise definitions. To overcome these biases by utilizing pre-trained word-embedding, authors learned embedding from existing input word embeddings using encoder-decoder architecture by defining decoder cost function that considers dictionary agreement as a constraint and decoded the debiased embedding~\cite{kaneko_dictionary_2021}. 

% \subsection{Methodology Comparisons}
% Mostly focus on different group of papers focused on similar task with similar
% datasets. show tables that mentioned - dataset names, algorithm used
% (highlevel), performances (acc, f1 or so on). Need to think after writing those
% previous sections.

% \subsection{Deep Learning}
% general diagram for deep learning approach.


% Semi-supervised approach \cite{patra_bilingual_2019}.

% \cite{wu_2021_sequence}We computationally model the processes of
% word borrowing from a donor word to an incorporated
% word, and vice versa, by answering
% two questions: (1) what does a word look
% like incorporated into another language, and
% in the opposite direction (2) where did a word
% come from? We experiment with several
% model variants, including LSTM encoderdecoders,
% copy attention, and Transformers.

% \cite{wu_computational_2020} For our model, we used a LSTM with an embedding
% dimension of 128 and hidden dimension of 128.
% The output of the last hidden state is passed to a fully connected layer with a sigmoid activation function.

% \subsubsection{Parameters and Evaluation}
% \begin{itemize}
%     \item Perplexity
%     \item Precision
%     \item BLEU score
%     \item ROUGE-L
%     \item Cosine similarity
% \end{itemize}

% In general, the performance of generative definition models are evaluated using
% perplexity and BLEU score.

\section{Available Datasets and Tasks}
This section discusses the available datasets used to solve different etymological tasks. The authors create, filter, or utilize existing data sources depending on the task. The most commonly used dataset is Wikitonary~\cite{wiktionary}, an online editable large multilingual online dictionary. Wiktionary dictionary contains the definition of words and related information to understand those definitions. Therefore, it includes etymologies, cognates, pronunciations, antonyms, synonyms, and translations. Researchers synthesize this dataset to create their task-specific dataset. Fourier~\cite{fourrier_methodological_2020} creates a lexical etymological database, EtymDB 2.0 from Wiktionary. Winston~\cite{wu_computational_2020} developed Wiktionary parser to parse etymology. Etymological Wordnet captures etymological word-formation information between two multilingual words~\cite{demelo_etymological_2014}. Starling is a derived dataset from the Uralic language family and contains over 2500 cognate sets\cite{nouri_alignment_2016}. Researchers proposed various etymological tasks such as parsing etymology, mapping cognates, establishing relationships between words to find word origin, modeling phylogenetic inferences that take aligned cognates, etc. We identify related etymological datasets and corresponding research tasks and summarize them in table~\ref{tab:datasets_task}.

\begin{longtable}
    {|p{2cm}|p{2cm}|p{2cm}|p{4cm}|p{4cm}|}
    \caption{Document Modeling Available Dataset and Task}
    \label{tab:datasets_task}                                                                                                                                                                                                                                                                                                                                                                                                                                                                                                                                                                                                                                                 \\\hline
    Article                                              & dataset                                                                                                           & types                                                                                         & Brief Description                                                                                                                                                                                                         & task                                                                                                                                                               \\ \hline
    Fourrier et. al. \cite{fourrier_methodological_2020} & EtymDB 2.0                                                                                                        & Lexical database                                                                              & This etymological dataset was generated form Wiktionary. Contains 1.8 million lexemes and 700,000 fine-grained etymological relations over                                                                                & Proposed guidelines to generate etimological dataset and also introduces usescases such as phylogenetic tree generation, low resource
    machine translation and medieval languages study.                                                                                                                                                                                                                                                                                                                                                                                                                                                                                                                                                                                                                         \\ \hline
    Gerard~\cite{demelo_etymological_2014}               & Etymological Wordnet                                                                                              & Definition                                                                                    & Contains multilingual, machine readable word origin networks of over  500,000 word origin created from Wiktionary                                                                                                         & Focuses on creating word-origin and linking using creation using custom pattern matching techniques.                                                               \\ \hline

    Benjamin \cite{mericli_annotating_2015}              & Turkic Language Dictionary                                                                                        & dictionary                                                                                    & This Turkic language dictionary contains 1996 English glosses and corresponding translation in eight Turkic languages such as  Turkish, Kyrgyz, Turkmen, Azerbaijani, Kazakh, Tatar,
    Uzbek, and Uyghur.                                   & The author proposed a methodology for annotating Cognates and Etymological Origin for Turkic language.                                                                                                                                                                                                                                                                                                                                                                                                                                                                                                             \\ \hline

    Jack (2016)~\cite{bowers_deep_2016}                  & -                                                                                                                 & generic                                                                                       & Digital dictionary                                                                                                                                                                                                        & present a set of modeling principles for representing etymological data using TEI framework.                                                                       \\ \hline


    Lisa (2020)~\cite{bowers_deep_2016}                  & Gathered from parrallel corpus                                                                                    & Multilingual                                                                                  & Word-based semantic analysis based on two major word-level corpus SWADESH Words, PEREIRA Words. Sentence level semantic analysis using corpus EUROPARL Sentences.                                                         & Proposed method for similarity analysis and identify semantic-drift between language families.                                                                     \\ \hline

    Barun et. al. (2019)~\cite{patra_bilingual_2019}     & MUSE and VecMap dataset                                                                                           & Multilingual                                                                                  & MUSE dataset contains embeddings on Wikipedia and bilingual dictionaries using Facebook internal tools. VecMap contains CBOW embedding on WacKy crawling corpora.                                                         & Quantitatively Estimate isometry between embedding.                                                                                                                \\ \hline


    Winston et.al. (2020)~\cite{wu_computational_2020}   & Wiktionary dictionary                                                                                             & Multilingual                                                                                  & online editable multilingual dictionary,                                                                                                                                                                                  & Developed parser for Wiktionary dataset and also predict etimology.                                                                                                \\ \hline

    Winston et.al. (2021)~\cite{wu_2021_sequence}\cite{} & Wiktionary (eng)                                                                                                  & English                                                                                       & etymology information from the English
    edition of Wiktionary using Yawipa.                  & Proposed computation methodology for establishing the directed etymology relationship between etymology between a
    donor word and an incorporated word.                                                                                                                                                                                                                                                                                                                                                                                                                                                                                                                                                                                                                                      \\ \hline

    Hannes (2012)~\cite{wettig_using_2012}               & StarLing, and The Origin of Finnish Words                                                                         & StarLing is Uralic Database. "The Origin of Finnish Words" is Finnish etymological dictionary & StarLing contains over 2500 cognates. The Finnish etymological dictionary has over 5000 cognate sets                                                                                                                      & Proposed methodology to find the best alignment at the sound level.                                                                                                \\ \hline

    Javad (2013)~\cite{nouri_alignment_2016}             & StarLing                                                                                                          & Uralic database                                                                               & Contains 2586 Uralic cognate sets                                                                                                                                                                                         & Used aligned cognates as input to model for phylogenetic inference
    from population genetics. In this task, they utilize sub-tree of Uralic, viz., the Finno-Ugric subfamily from Uralic database.                                                                                                                                                                                                                                                                                                                                                                                                                                                                                                                                            \\ \hline

    Carlo (2019)\cite{DBLP:journals_corr_abs-1908-02477} & Ab Antiquo                                                                                                        & Multilingual                                                                                  & Dataset contains cognates and their proto-words in both orthographic and phonetic forms in six different languages. It has in total of 8,799 cognate sets.                                                                & Augment dataset from Wiktionary and developed comparative methods that takes daughter language cognates and identify the form of proto-words in ancestor language. \\ \hline

    Lisa (2013)\cite{beinborn-etal-2013-cognate}         & multilingual resources UBY and Universal WordNet (UWN)                                                            & Multilingual Word pairs                                                                       & Over 1000 cognate pairs.                                                                                                                                                                                                  & Generate cognates using character-based machine translation  approach and evaluate model on three language
    pairs English-Russian, English-Greek, and
    English-Farsi.                                                                                                                                                                                                                                                                                                                                                                                                                                                                                                                                                                                                                                                            \\ \hline

    Sagot (2017)\cite{Sagot2017ExtractingAE}             & Etymological Dataset                                                                                              & Multilingual                                                                                  & etymological dataset containing nearly one million lexemes and half million relationship between lexemes.                                                                                                                 & Extract etymology from Wiktionary.                                                                                                                                 \\\hline

    Vivi (2013)\cite{nastase-strapparava-2013-bridging}  & Comparable Corpora of News Article                                                                                & Multilingual                                                                                  & This dataset consists of comparable corpora of news articles in English and Italian. Each news article is annotated with any of given four categories. It contains in total of 27821 English and 32354 Italian documents. & Proposed methodology for cross-language document categorization using etymological ancestors of the words of the documents.                                        \\ \hline

    % \end{tabular}
\end{longtable}

% \begin{table}[htbp]
%     \centering
%     \begin{tabular}{|p{2cm}|p{2cm}|p{2cm}|p{4cm}|p{4cm}|}\hline
%     Article & dataset & types & Brief Description & task\\ \hline
%         Fourrier et. al. \cite{fourrier_methodological_2020} & EtymDB 2.0 & Lexical database & This etymological dataset was generated form Wiktionary. Contains 1.8 million lexemes and 700,000 fine-grained etymological relations over  &  Proposed guidelines to generate etimological dataset and also introduces usescases such as phylogenetic tree generation, low resource
% machine translation and medieval languages study. \\ \hline
% Gerard \cite{demelo_etymological_2014}& Etymological Wordnet & Definition & Contains multilingual, machine readable word origin networks of over  500,000 word origin created from Wiktionary & Focuses on creating word-origin and linking using creation using custom pattern matching techniques. \\ \hline

% Benjamin \cite{mericli_annotating_2015}& Turkic Language Dictionary & dictionary & This Turkic language dictionary contains 1996 English glosses and corresponding translation in eight Turkic languages such as  Turkish, Kyrgyz, Turkmen, Azerbaijani, Kazakh, Tatar, 
%  Uzbek, and Uyghur. & The author proposed a methodology for annotating Cognates and Etymological Origin for Turkic language.\\ \hline

%  Jack (2016)~\cite{bowers_deep_2016} & - & generic & Digital dictionary & present a set of modeling principles for representing etymological data using TEI framework. \\ \hline


%   Lisa (2020)~\cite{bowers_deep_2016} & Gathered from parrallel corpus& Multilingual & Word-based semantic analysis based on two major word-level corpus SWADESH Words, PEREIRA Words. Sentence level semantic analysis using corpus EUROPARL Sentences.  & Proposed method for similarity analysis and identify semantic-drift between language families.\\ \hline

%   Barun et. al. (2019)~\cite{patra_bilingual_2019} & MUSE and VecMap dataset & Multilingual & MUSE dataset contains embeddings on Wikipedia and bilingual dictionaries using Facebook internal tools. VecMap contains CBOW embedding on WacKy crawling corpora. & Quantitatively Estimate isometry between embedding. \\ \hline


%     Winston et.al. (2020)~\cite{wu_computational_2020} & Wiktionary dictionary & Multilingual  & online editable multilingual dictionary,  & Developed parser for Wiktionary dataset and also predict etimology. \\ \hline

%     Winston et.al. (2021)~\cite{wu_2021_sequence}\cite{} & Wiktionary (eng) & English & etymology information from the English
% edition of Wiktionary using Yawipa. & Proposed computation methodology for establishing the directed etymology relationship between etymology between a
% donor word and an incorporated word.  \\ \hline

%     Hannes (2012)~\cite{wettig_using_2012} & StarLing, and The Origin of Finnish Words & StarLing is Uralic Database. "The Origin of Finnish Words" is Finnish etymological dictionary& StarLing contains over 2500 cognates. The Finnish etymological dictionary has over 5000 cognate sets & Proposed methodology to find the best alignment at the sound level. \\ \hline

%     Javad (2013)\cite{nouri_alignment_2016} & StarLing  & Uralic database & Contains 2586 Uralic cognate sets & Used aligned cognates as input to model for phylogenetic inference
% from population genetics. In this task, they utilize sub-tree of Uralic, viz., the Finno-Ugric subfamily from Uralic database. \\ \hline

%     Carlo (2019)\cite{DBLP:journals_corr_abs-1908-02477} & Ab Antiquo & Multilingual& Dataset contains cognates and their proto-words in both orthographic and phonetic forms in six different languages. It has in total of 8,799 cognate sets.  & Augment dataset from Wiktionary and developed comparative methods that takes daughter language cognates and identify the form of proto-words in ancestor language.   \\ \hline

%     author (year)\cite{} & dataset & type& data descp & task \\ \hline

%     \end{tabular}
%     \caption{Document Modeling Available Dataset and Task}
%     \label{tab:datasets_task}
% \end{table}
\FloatBarrier
\section{Challenges and future directions}
Definition modeling faces several challenges, providing new opportunities for future research to be developed.

\noindent\textbf{Polysemes:} The basic definition model cannot be used to generate definitions for \textit{polysemes}, words with multiple definitions. As a
significant challenge for the original definition model, many researchers have
proposed methods to tackle this problem. However, many of the proposed
approaches require the context of the \textit{definiendum} to be provided to the
model. Methods that provide appropriate definitions for polysemes without
context may be valuable in tasks with limited language resources.

\noindent\textbf{Technical terms:} It is challenging to generate definitions for technical terms which require expert knowledge of the field \cite{huang_cdm_2021}. It may
be necessary to provide specific context to generate
definitions for technical terms appropriately. However, obtaining the context requires
scraping and parsing web resources outside of the standard datasets available. Therefore, it may be necessary to generate definitions for technical
terms to augment dictionary datasets properly.

\noindent\textbf{Word combinations:} Complex word combinations, including proverbs and sayings, are rarely covered by sense inventories~\cite{bevilacqua_generationary_2020}. In word combinations, multiple words are used in series to create a new phrase that may be interpreted as a single word for the case of definition modeling. Since the resulting definition
of word combinations may or may not depend on the words used, context may be necessary to parse these word combinations and generate useful definitions. Still, more research is needed to determine if this is the case.
Additionally, word combinations may be absent from the standard dictionary-based datasets.

\noindent\textbf{Non-English words:} As many of the datasets developed for defintion modeling thus far take information from English dictionaries, most methods also
are only applied to English words. In addition, as there exist several lexical resources in other languages, it should be possible to generate definitions for non-English words. To evaluate the quality of word embeddings for
non-English words within definition modeling, it is necessary to develop a method to generate definitions for non-English words. There is some work in Chinese definition modeling \cite{zheng_decompose_2021}, and in French
definition modeling~\cite{reid_vcdm_2020}. However, more research is needed to determine the best method for generating definitions for non-English words, 
especially for a model that can generalize across multiple languages.

\noindent\textbf{Evaluation criteria:} Definition models have been evaluated on a number of metrics, including precision, perplexity, BLEU, and ROUGE. However,
as definition modeling aims to improve the interpretability of word embeddings, it is important to select the evaluation criteria correctly. Many definitions consist of a single word, which can interfere with evaluation metrics such as BLEU and ROUGE scores \cite{mickus_mark_2019}. Human
evaluation of generated definitions can be useful but difficult for researchers to obtain.


\FloatBarrier
\section{Conclusion}
Write text related to conclusion of this study.

\printbibliography
\end{document}
